\section{Logarithmic squares, Gamma integrals, and a dilogarithmic primitive}
\label{hqj:sec:gamma}

The simplest nonlocal defect has a useful independent calculus.
Let $a>0$, and let $c,d\in\C$ satisfy
$\Rea(a+c)>0$ and $\Rea(a+d)>0$. Put $h=c-d$ and
\begin{equation}\label{hqj:eq:Sdef}
\mathscr S(a;c,d)=\sum_{n\ge0}
\bigl(\log(n+a+c)-\log(n+a+d)\bigr)^2.
\end{equation}
The summand is $O(n^{-2})$, uniformly on compact sets of these
parameters, so this is an ordinary holomorphic sum. In particular
$\mathscr S(a;ib,-ib)=\Delta_2(1;a,b)$.

\begin{theorem}[A digamma logarithmic integral]\label{hqj:thm:gamma-square}
On the domain above,
\begin{equation}\label{hqj:eq:Sdigamma}
\boxed{\mathscr S(a;c,d)=2h\int_0^1
\psi(a+d+ht)\log\frac{t}{1-t}\dd t.}
\end{equation}
The integral is ordinary and convergent at both endpoints. Moreover,
\begin{equation}\label{hqj:eq:Smixed}
\partial_c\partial_d\mathscr S(a;c,d)
=-2\frac{\psi(a+d)-\psi(a+c)}{d-c},
\end{equation}
with its removable value $-2\psi'(a+c)$ when $c=d$.
Every derivative in $a$ is obtained by replacing $\psi$ by the
corresponding polygamma in \eqref{hqj:eq:Sdigamma}.
\end{theorem}
\begin{proof}
For each $n$, integrate the resolvent along the segment from $d$ to $c$:
\[
\log(n+a+c)-\log(n+a+d)
=h\int_0^1\frac{\dd t}{n+a+d+ht}.
\]
The entire segment lies in the right half-plane after adding $a$.
Compact uniform resolvent estimates justify squaring, summing, and
interchanging the two parameter integrals. The standard digamma
difference identity gives
\[
\sum_{n\ge0}\frac1{(n+A)(n+B)}
=\frac{\psi(B)-\psi(A)}{B-A}.
\]
Thus $\mathscr S$ equals
\[
h\int_0^1\int_0^1
\frac{f(v)-f(t)}{v-t}\dd v\dd t,
\qquad f(t)=\psi(a+d+ht).
\]
The divided difference has a removable diagonal and is continuous.
For a continuously differentiable $f$, symmetric exclusion of a thin
diagonal strip gives
\[
\int_0^1\int_0^1\frac{f(v)-f(t)}{v-t}\dd v\dd t
=2\int_0^1f(t)\log\frac{t}{1-t}\dd t.
\]
This follows by interchanging the two terms on the symmetric truncated
domain and then integrating $1/(t-v)$; the endpoint logarithms are
integrable and the omitted divided-difference strip tends to zero.
This proves \eqref{hqj:eq:Sdigamma}, including complex parameters.

Differentiating the normally convergent sum in $c,d$ gives
$-2\sum((n+a+c)(n+a+d))^{-1}$, which is
\eqref{hqj:eq:Smixed}. Differentiation under the integral for a fixed
number of $a$ derivatives is justified by the same right-half-plane
compactness and integrable endpoint logarithms.
\end{proof}

\subsection{An anchored primitive with explicit dilogarithms}
For the next representation assume real $c>d$ and $a+d>0$, and set $h=c-d>0$.
For $y>0$ define
\begin{equation}\label{hqj:eq:fprimitive}
f_h(y)=y\log^2(1+h/y)-2h\Li_2\!\left(\frac{h}{y+h}\right).
\end{equation}
Direct differentiation, using the ordinary dilogarithm derivative,
gives
\begin{equation}\label{hqj:eq:fderivative}
f_h'(y)=\log^2(1+h/y),\qquad
f_h(y)=-\frac{h^2}{y}+O(y^{-2})\quad(y\longrightarrow\infty).
\end{equation}
The counterterm in the next theorem is therefore necessary.

\begin{theorem}[Dilogarithmic and Gamma representations of the primitive]
\label{hqj:thm:primitive}
The series
\begin{equation}\label{hqj:eq:Pprimitive}
\mathscr P(a;c,d)=h^2\psi(a+d)+
\sum_{n\ge0}\left\{f_h(n+a+d)+\frac{h^2}{n+a+d}\right\}
\end{equation}
is normally convergent for real $a+d>0$ on compact subintervals, and
\begin{equation}\label{hqj:eq:Pderivative}
\frac{\dd}{\dd a}\mathscr P(a;c,d)=\mathscr S(a;c,d).
\end{equation}
For $a_0+d,a_1+d>0$,
\begin{align}
\mathscr P(a_1;c,d)-\mathscr P(a_0;c,d)
={}&2h\int_0^1
\log\frac{\Gamma(a_1+d+ht)}{\Gamma(a_0+d+ht)}
\log\frac{t}{1-t}\dd t.
\label{hqj:eq:gamma-dilog}
\end{align}
Thus \eqref{hqj:eq:gamma-dilog} is an exact identity between an ordinary
log-Gamma integral and an explicitly convergent dilogarithmic sum.
\end{theorem}
\begin{proof}
The asymptotic formula in \eqref{hqj:eq:fderivative} leaves an
$O(n^{-2})$ remainder in \eqref{hqj:eq:Pprimitive}, and differentiating
that remainder gives $O(n^{-3})$ on compact sets. Its derivative is
\[
h^2\psi'(a+d)+\sum_{n\ge0}
\left\{\log^2(1+h/(n+a+d))-\frac{h^2}{(n+a+d)^2}\right\}.
\]
Since $\psi'(z)=\sum_{n\ge0}(n+z)^{-2}$, the rational terms cancel
and give \eqref{hqj:eq:Pderivative}. Integrate
\eqref{hqj:eq:Sdigamma} between $a_0$ and $a_1$. Fubini is valid on the
compact $a$ interval, and the primitive of $\psi$ is the real
$\log\Gamma$. This gives \eqref{hqj:eq:gamma-dilog} with the displayed
anchoring, not an unspecified additive constant.
\end{proof}

For $|h|<a+d$, the same family has an independent harmonic expansion
\begin{equation}\label{hqj:eq:S-harmonic}
\mathscr S(a;d+h,d)=
2\sum_{r\ge2}\frac{(-1)^rH_{r-1}}r h^r\zeta(r,a+d).
\end{equation}
This follows from
$\log^2(1+z)=2\sum_{r\ge2}(-1)^rH_{r-1}z^r/r$.
Its anchored $a$ primitive uses
\[
\int_{a_0}^{a_1}\zeta(r,a+d)\dd a
=\begin{cases}
\psi(a_1+d)-\psi(a_0+d),&r=2,\\
\dfrac{\zeta(r-1,a_0+d)-\zeta(r-1,a_1+d)}{r-1},&r>2.
\end{cases}
\]
This gives another exact bridge between the Gamma integral,
harmonic coefficients, and Hurwitz values, and supplies an independent
numerical route used in the accompanying tests.
